\section{Recommendation}
\label{sec:recommendation}

\subsection{Action 1: fix the promises (largest effect)}
The cheapest intervention is to change the promise. \Cref{tab:scenarios} applies cumulative
changes to the test-period orders, using the noon cut-off learned on the training period.

\begin{table}[!htbp]
  \centering
  \caption{Promise scenarios on the test period (428 orders per week). Late cost is $0.10\times$ order value per late order.}
  \label{tab:scenarios}
  \footnotesize
  \setlength{\tabcolsep}{4pt}
  \begin{tabular}{l c c c c}
    \toprule
    \textbf{Scenario} & \textbf{Late share} & \textbf{Late/wk} & \textbf{Mean promise} & \textbf{Late cost/wk} \\
    \midrule
    Current (0/1/2/4 days)                          & 57.6\% & 246 & 2.93 d & \usd{10{,}520} \\
    A: Same Day after 12:00 promised next day       & 54.5\% & 233 & 2.96 d & \usd{10{,}020} \\
    B: A + First Class promised 2 days              & 39.4\% & 169 & 3.11 d & \usd{7{,}240} \\
    C: B + Second Class promised 4 days             & 31.5\% & 135 & 3.50 d & \usd{5{,}770} \\
    D: C + Second/Standard promised 6 days          & 0.0\%  & 0   & 5.09 d & \usd{0} \\
    \bottomrule
  \end{tabular}
\end{table}

\textbf{Recommended: scenario B, now.} First Class already takes exactly 2 days, so the new
promise costs no speed, and a checkout cut-off for Same Day is standard practice. B removes 77
late orders and about \usd{3{,}300} of late cost per week. Scenario C is a commercial decision:
Second Class is no faster than Standard, so DataCo should fix it or stop selling it as faster.
We do not recommend D, whose lost conversions the data cannot measure.

\subsection{Action 2: flag orders at release with the lookup table}
Some late orders remain after the promise change, so the planner still needs a daily list.
\Cref{fig:saving} plots weekly net saving against the share of orders flagged, with orders
ranked by $\hat p\times v$. The saving peaks at about a 45\% alert rate for both the lookup
table and the decision tree (\usd{2{,}400} per week). Beyond that point, each extra check costs
more than the late delivery it prevents. Rule B peaks earlier, at 25\% (\usd{1{,}810} per week),
because it cannot tell a Second Class order (80\% late) from a First Class one (100\% late).

\begin{figure}[!htbp]
  \centering
  \includegraphics[width=0.66\linewidth]{05_saving_vs_alert_rate}
  \caption{Weekly net saving compared with taking no action, by alert rate, on the test period. The dotted line marks Rule A's alert rate.}
  \label{fig:saving}
\end{figure}

\begin{table}[!htbp]
  \centering
  \caption{Weekly operational effect of each policy on the test period (428 orders per week).}
  \label{tab:rec}
  \footnotesize
  \setlength{\tabcolsep}{3pt}
  \begin{tabular}{l c c c c c}
    \toprule
    \textbf{Policy} & \textbf{Flagged} & \textbf{Late caught} & \textbf{Late missed} & \textbf{False alarms} & \textbf{Net saving} \\
    \midrule
    Lookup table, EV rule & 187 & 129 & 117 & 57  & \usd{2{,}399} \\
    Decision tree, EV rule & 188 & 129 & 117 & 59  & \usd{2{,}393} \\
    Rule B (+ Same Day cut-off)        & 162 & 146 & 101 & 16  & \usd{1{,}503} \\
    Rule A (First/Second)              & 149 & 133 & 113 & 16  & \usd{1{,}381} \\
    Flag every order                   & 428 & 246 & 0   & 181 & \usd{986} \\
    \bottomrule
  \end{tabular}
\end{table}

\textbf{Recommended policy:} each morning, flag the orders for which
$P(\text{late})\times\text{order value} > \usd{200}$, where $P(\text{late})$ comes from the
five-row lookup table below. The expected workload is about \textbf{190 orders per week},
roughly 27 per day. That is about 16 hours of customer-service time per week, assuming 5
minutes per order. It catches 129 late orders per week and saves about \usd{900} per week more
than Rule B.

\begin{center}
  \small
  \begin{tabular}{l c c}
    \toprule
    \textbf{Cell (from training data)} & \textbf{$P(\text{late})$} & \textbf{Flag if order value exceeds} \\
    \midrule
    First Class; Same Day after 12:00 & 1.00 & \usd{200} \\
    Second Class                      & 0.80 & \usd{250} \\
    Standard Class                    & 0.40 & \usd{500} \\
    Same Day before 12:00             & 0.00 & never \\
    \bottomrule
  \end{tabular}
\end{center}

\textbf{Why the lookup table and not the model.} The tree and gradient boosting give the same
answer at extra cost: retraining, monitoring, and logic planners cannot check by eye. Across all
36 cost settings (action \usd{2}--\usd{20}, miss 5--20\% of value, 30--70\% avoided), the best model
and the table differ by $-\usd{24}$ to $+\usd{5}$ per week, while the table beats Rule B by
\usd{300}--\usd{5{,}500} per week. At a \usd{2} action cost the optimal alert rate rises to 60--90\%.

\paragraph{Size of the effect.} Under the stated assumptions, scenario B avoids about
\textbf{\usd{3{,}300} per week} of late-delivery cost, and the flagging policy saves a further
\textbf{\usd{900} per week} over the current-practice proxy: roughly \usd{200}k per year at 2017
volumes. Because the saving scales with order value, which fell by two-thirds in late 2017, we
would quote it as ``tens to low hundreds of thousands of dollars per year''.
